% English translation of questions/5780/semA-moedB/q4.tex (metadata is taken from the Hebrew file).

\begin{question}
Let $f(x)$ be a function such that $f(x)+\sin(x)$ is differentiable on the open interval $(a,b)$. Which of the following must be true?
\begin{enumerate}
  \item $f(x)+\cos(x)$ is differentiable on the open interval $(a,b)$.
  \item $f(x)+\sin(x)$ is continuous on the closed interval $[a,b]$.
  \item $f(x)$ is continuous on the closed interval $[a,b]$.
  \item $f(x)$ is bounded on the open interval $(a,b)$.
  \item $f(x)+\sin(x)$ is bounded on the open interval $(a,b)$.
\end{enumerate}
\end{question}

\begin{hint}
Write $f+\cos$ as the sum of $f+\sin$ and a known differentiable function.
\end{hint}

\begin{solution}
\textbf{The correct answer: a.}

Denote $g(x)=f(x)+\sin x$, differentiable on $(a,b)$. Then
\[ f(x)+\cos x=g(x)-\sin x+\cos x , \]
and this is a sum of functions differentiable on $(a,b)$ ($\sin$ and $\cos$ are differentiable on all of $\R$), hence it is differentiable on $(a,b)$. That is, (a) is always true.

The other statements need not hold. Counterexample: $f(x)=\frac{1}{x-a}-\sin x$ on $(a,b)$. Then $f(x)+\sin x=\frac{1}{x-a}$ is differentiable on $(a,b)$, but:
\begin{itemize}
  \item $f+\sin$ and $f$ are not defined (and certainly not continuous) at $x=a$ --- (b) and (c) fail. (Even if we assign them some value at $a$, they are unbounded near $a$ and hence not continuous there.)
  \item $\frac{1}{x-a}\to+\infty$ as $x\to a^+$, hence $f+\sin$ is unbounded on $(a,b)$ --- (e) fails; and also $f=\frac{1}{x-a}-\sin x$ is unbounded (since $\sin$ is bounded) --- (d) fails.
\end{itemize}
\end{solution}
